% GENERATED FILE — DO NOT EDIT. Written by 01_code/04_tables_values/make_tables_values.py
% and copied to manuscript/values.tex by step 3/5 of run_all.sh. Every edit
% made here is overwritten at the next run; change the generator instead. The trailing
% comment on each macro is REPORTED, not authored: it is the 'source' string
% that collect_computed.py attached to the value, and it describes where the
% number came from, not what the macro means. This is not a source file.
%
% values_generated.tex — CANDIDATE values.tex written by make_tables_values.py.
% DO NOT EDIT BY HAND: regenerate with run_all.sh (or make_tables_values.py).
% Drop-in for \input{values.tex}: every macro is defined.
% Block 1 = sourced literature inputs; Block 2 = code-regenerated outputs.
% There is no third block. Every macro here is either sourced or computed;
% a macro that is neither fails the run by name rather than being copied
% through from the previous run's output.
% Generation date = the commit that carries this file; no timestamp is written
% inline, so that two runs of the same inputs stay byte-identical.

% ===== (1) Literature inputs (sourced constants, NOT computed) =====
\newcommand{\DTstarAMOCHi}{8.0}   % Armstrong McKay et al. 2022, Science — Table 1, AMOC collapse, high end of the reported threshold range (°C)
\newcommand{\DTstarAMOCLo}{1.4}   % Armstrong McKay et al. 2022, Science — Table 1, AMOC collapse, low end of the reported threshold range (°C)
\newcommand{\DTstarAMOC}{4}   % Armstrong McKay et al. 2022, Science — central threshold
\newcommand{\DTstarAmazonHi}{6.0}   % Armstrong McKay et al. 2022, Science — Table 1, Amazon dieback, high end of the reported threshold range (°C)
\newcommand{\DTstarAmazonLo}{2.0}   % Armstrong McKay et al. 2022, Science — Table 1, Amazon dieback, low end of the reported threshold range (°C)
\newcommand{\DTstarAmazon}{3.5}   % Armstrong McKay et al. 2022, Science — central threshold
\newcommand{\DTstarWAISHi}{3.0}   % Armstrong McKay et al. 2022, Science — Table 1, West Antarctic ice sheet, high end of the reported threshold range (°C)
\newcommand{\DTstarWAISLo}{1.0}   % Armstrong McKay et al. 2022, Science — Table 1, West Antarctic ice sheet, low end of the reported threshold range (°C)
\newcommand{\DTstarWAIS}{1.5}   % Armstrong McKay et al. 2022, Science — central threshold
\newcommand{\EglobalRate}{40}   % Global Carbon Project — gross global emissions ~40 GtCO2/yr
\newcommand{\Lmax}{0.67}   % Dietz et al. 2021, PNAS — upper end of the common damage sweep
\newcommand{\NMonteCarlo}{100{,}000}   % Monte Carlo draws in calibration_ci.py, eps_inf_ci.py and fig_forest_tc.py (modelling choice)
\newcommand{\TAMOCHi}{300}   % Armstrong McKay et al. 2022, Science — Table 1, AMOC collapse, timescale, high end of range (yr)
\newcommand{\TAMOCLo}{15}   % Armstrong McKay et al. 2022, Science — Table 1, AMOC collapse, timescale, low end of range (yr)
\newcommand{\TAMOC}{50}   % Armstrong McKay et al. 2022, Science — Table 1 central
\newcommand{\TAmazonHi}{200}   % Armstrong McKay et al. 2022, Science — Table 1, Amazon dieback, timescale, high end of range (yr)
\newcommand{\TAmazonLo}{50}   % Armstrong McKay et al. 2022, Science — Table 1, Amazon dieback, timescale, low end of range (yr)
\newcommand{\TAmazon}{100}   % Armstrong McKay et al. 2022, Science — Table 1 central
\newcommand{\TCRElikelyHi}{0.63}   % IPCC AR6 WG1 2021, SPM D.1.1, likely range — upper end, °C per 1000 GtCO2
\newcommand{\TCRElikelyLo}{0.27}   % IPCC AR6 WG1 2021, SPM D.1.1, likely range — lower end, °C per 1000 GtCO2
\newcommand{\TCREmean}{0.45}   % IPCC AR6 WG1 2021 — °C per 1000 GtCO2 (central)
\newcommand{\TCREsd}{0.19}   % computed: (TCRE_likely_hi - TCRE_likely_lo) / 2 / z(0.83), the AR6 likely range (SPM D.1.1) as the 17th-83rd centiles of a normal law, °C per 1000 GtCO2 (0.18865)
\newcommand{\TWAISHi}{13{,}000}   % Armstrong McKay et al. 2022, Science — Table 1, West Antarctic ice sheet, timescale, high end of range (yr)
\newcommand{\TWAISLo}{500}   % Armstrong McKay et al. 2022, Science — Table 1, West Antarctic ice sheet, timescale, low end of range (yr)
\newcommand{\TWAIS}{2000}   % Armstrong McKay et al. 2022, Science — Table 1 central
\newcommand{\bhpChiPost}{2.4}   % Hambel, van der Ploeg and van den Bremer, replication package, Zenodo 20075290 (2026-05-07), Code_Parameters.m / Code_Initiation.m, as printed by MakeTable1.m — chi_post, TCRE post-tip, degC/TtC
\newcommand{\bhpChiPre}{1.8}   % Hambel, van der Ploeg and van den Bremer, replication package, Zenodo 20075290 (2026-05-07), Code_Parameters.m / Code_Initiation.m, as printed by MakeTable1.m — chi_pre, TCRE pre-tip, degC/TtC
\newcommand{\bhpHone}{0.35\%\,{}^{\circ}\mathrm{C}^{-1}\,\mathrm{yr}^{-1}}   % Hambel, van der Ploeg and van den Bremer, replication package, Zenodo 20075290 (2026-05-07), Code_Parameters.m / Code_Initiation.m, as printed by MakeTable1.m — h_1T, marginal arrival rate of the first tipping point, %/degC/yr
\newcommand{\bhpRungs}{5}   % Hambel, van der Ploeg and van den Bremer, replication package, Zenodo 20075290 (2026-05-07), Code_Parameters.m / Code_Initiation.m, as printed by MakeTable1.m — transitions between the H states
\newcommand{\bhpStates}{6}   % Hambel, van der Ploeg and van den Bremer, replication package, Zenodo 20075290 (2026-05-07), Code_Parameters.m / Code_Initiation.m, as printed by MakeTable1.m — H, number of tipping states
\newcommand{\bhpTstar}{1}   % Hambel, van der Ploeg and van den Bremer, replication package, Zenodo 20075290 (2026-05-07), Code_Parameters.m / Code_Initiation.m, as printed by MakeTable1.m — T*, tipping threshold, degC
\newcommand{\bhpTzero}{1.1}   % Hambel, van der Ploeg and van den Bremer, replication package, Zenodo 20075290 (2026-05-07), Code_Parameters.m / Code_Initiation.m, as printed by MakeTable1.m — T_0, initial temperature anomaly, degC
\newcommand{\dietzAmazonSCCadd}{0.04}   % Dietz, Rising, Stoerk and Wagner (2021), PNAS 118(34), e2103081118; value read and checked by the author; Table 2, expected SCC with Amazon dieback (52.07) minus 'None' (52.03), 2020 US$/tCO2; the table prints the two levels, not the difference
\newcommand{\dietzAmazonSCCpct}{0.1}   % Dietz, Rising, Stoerk and Wagner (2021), PNAS 118(34), e2103081118; value read and checked by the author; Table 2, 'Increase due to TP, %', Amazon; text: 'Dieback of the Amazon rainforest leads to a modest 0.1% increase in the expected SCC'
\newcommand{\dietzSCCbase}{52.03}   % Dietz, Rising, Stoerk and Wagner (2021), PNAS 118(34), e2103081118; value read and checked by the author; Table 2, expected SCC with no tipping point ('None'), 2020 US$/tCO2
\newcommand{\relaxWAISHillHi}{20}   % Hill et al. 2023, The Cryosphere 17:3739-3759, doi:10.5194/tc-17-3739-2023 — e-folding relaxation time of the grounding-line flux after a 20-yr ocean perturbation, Antarctic-wide, high end of 'between 10 and 20 years' (yr)
\newcommand{\relaxWAISHillLo}{10}   % Hill et al. 2023, The Cryosphere 17:3739-3759, doi:10.5194/tc-17-3739-2023 — e-folding relaxation time of the grounding-line flux after a 20-yr ocean perturbation, Antarctic-wide, low end of 'between 10 and 20 years' (yr)
\newcommand{\rhoBase}{0.03}   % modelling choice — baseline discount rate (Nordhaus–Stern midpoint)
\newcommand{\sigmaPriorHi}{0.45}   % upper end of the common log-uniform noise prior of the Amazon and the ice sheet (modelling choice, common to the two elements)
\newcommand{\sigmaPriorLo}{0.05}   % lower end of the common log-uniform noise prior of the Amazon and the ice sheet (modelling choice, common to the two elements)
\newcommand{\sigmaPriorMed}{0.15}   % closed form: quantile 0.5 of the log-uniform prior on [0.05, 0.45]
\newcommand{\sigmaPriorPninety}{0.36}   % closed form: quantile 0.9 of the log-uniform prior on [0.05, 0.45]
\newcommand{\sigmaPriorPten}{0.06}   % closed form: quantile 0.1 of the log-uniform prior on [0.05, 0.45]

% ===== (2) Computed outputs (regenerated by the code) =====
\newcommand{\AMOCrecordFirstYear}{1871}   % ews_calibration.py — first year of the series, read from the header of 00_data/amoc/sg_index_hadisst.txt  [requires network]
\newcommand{\AMOCrecordLastYear}{2016}   % ews_calibration.py — last year of the series, read from the header of 00_data/amoc/sg_index_hadisst.txt  [requires network]
\newcommand{\DTzero}{1.2}   % Forster et al. 2024, Indicators of Global Climate Change 2023, ESSD 16:2625-2658 — observed warming above pre-industrial, 2014-2023 average, 1.19 [1.06 to 1.30] (°C)
\newcommand{\EglobalRateSSPhigh}{48.0}   % fig_forest_tc.py via emission_paths.py — SSP3-7.0 CO2 emissions in T0, GtCO2/yr, linear between the reported years; RCMIP v5.1.0 (Zenodo 4589756), Emissions|CO2, World, with the extensions of Meinshausen et al. 2020; extract in 00_data/emissions/
\newcommand{\EglobalRateSSPmid}{41.8}   % fig_forest_tc.py via emission_paths.py — SSP2-4.5 CO2 emissions in T0, GtCO2/yr, linear between the reported years; RCMIP v5.1.0 (Zenodo 4589756), Emissions|CO2, World, with the extensions of Meinshausen et al. 2020; extract in 00_data/emissions/
\newcommand{\LmaxFlowPct}{2}   % closed form: 100 rho Lmax, rho=0.03 (rhoBase), Lmax=0.67 (Lmax), the flow loss at the top of the damage range (%)
\newcommand{\MstarAMOCmed}{9433}   % calibration_ci.py — M*_P50 (MC seed=42)
\newcommand{\MstarAmazonmed}{7998}   % calibration_ci.py — M*_P50 (MC seed=42)
\newcommand{\MstarWAISmed}{3895}   % calibration_ci.py — M*_P50 (MC seed=42)
\newcommand{\OmegaAmazonmed}{0.32}   % calibration_ci.py — Omega(rho H) median, the realization factor
\newcommand{\OmegaWAISmed}{0.017}   % calibration_ci.py — Omega(rho H) median, the realization factor
\newcommand{\TCREdiscardPct}{0.9}   % calibration_ci.py — tcre_draw.draw_tcre, share of the normal draws <= 0 discarded by rejection, % of all draws made (seed=42)
\newcommand{\VarAMOChat}{0.088}   % ews_calibration.py — detrended variance (HadISST)  [requires network]
\newcommand{\ampAMOCSSPhigh}{4.0}   % closed form: F = (eps_now/eps_dyn)^(1/2) if eps_dyn < eps_now, else 1; the factor by which the leading law lifts the price between today's proximity and eps_dyn, the width of the window (the lift of the full price only if the state is inside the noise layer today); closed form, self-consistent: ((E/(M* theta)) sqrt(eps_now))^(2/3), E=82.7 GtCO2/yr (SSP3-7.0, highest rate before the median date, fig_forest_tc.py), M*=9433 GtCO2 (calibration_ci.py), theta=0.7928/yr measured on the record (ews_calibration.py), eps_now=0.7; theta declines as sqrt(eps) along the descent  [requires network]
\newcommand{\ampAMOCSSPmid}{4.9}   % closed form: F = (eps_now/eps_dyn)^(1/2) if eps_dyn < eps_now, else 1; the factor by which the leading law lifts the price between today's proximity and eps_dyn, the width of the window (the lift of the full price only if the state is inside the noise layer today); closed form, self-consistent: ((E/(M* theta)) sqrt(eps_now))^(2/3), E=44.3 GtCO2/yr (SSP2-4.5, highest rate before the median date, fig_forest_tc.py), M*=9433 GtCO2 (calibration_ci.py), theta=0.7928/yr measured on the record (ews_calibration.py), eps_now=0.7; theta declines as sqrt(eps) along the descent  [requires network]
\newcommand{\ampAMOC}{5.1}   % closed form: F = (eps_now/eps_dyn)^(1/2) if eps_dyn < eps_now, else 1; the factor by which the leading law lifts the price between today's proximity and eps_dyn, the width of the window (the lift of the full price only if the state is inside the noise layer today); closed form, self-consistent: ((E/(M* theta)) sqrt(eps_now))^(2/3), E=40 GtCO2/yr (E_0), M*=9433 GtCO2 (calibration_ci.py), theta=0.7928/yr measured on the record (ews_calibration.py), eps_now=0.7; theta declines as sqrt(eps) along the descent  [requires network]
\newcommand{\ampAmazonBoultonSSPhigh}{8.5}   % closed form: F = (eps_now/eps_dyn)^(1/2) if eps_dyn < eps_now, else 1; the factor by which the leading law lifts the price between today's proximity and eps_dyn, the width of the window (the lift of the full price only if the state is inside the noise layer today); closed form, self-consistent: ((E/(M* theta)) sqrt(eps_now))^(2/3), E=82.7 GtCO2/yr (SSP3-7.0, highest rate before the median date, fig_forest_tc.py), M*=7998 GtCO2 (calibration_ci.py), theta=9.677/yr, the vegetation-index time scale of Boulton et al. 2022, eps_now=0.66; theta declines as sqrt(eps) along the descent
\newcommand{\ampAmazonBoultonSSPmid}{10.5}   % closed form: F = (eps_now/eps_dyn)^(1/2) if eps_dyn < eps_now, else 1; the factor by which the leading law lifts the price between today's proximity and eps_dyn, the width of the window (the lift of the full price only if the state is inside the noise layer today); closed form, self-consistent: ((E/(M* theta)) sqrt(eps_now))^(2/3), E=44.3 GtCO2/yr (SSP2-4.5, highest rate before the median date, fig_forest_tc.py), M*=7998 GtCO2 (calibration_ci.py), theta=9.677/yr, the vegetation-index time scale of Boulton et al. 2022, eps_now=0.66; theta declines as sqrt(eps) along the descent
\newcommand{\ampAmazonBoulton}{10.8}   % closed form: F = (eps_now/eps_dyn)^(1/2) if eps_dyn < eps_now, else 1; the factor by which the leading law lifts the price between today's proximity and eps_dyn, the width of the window (the lift of the full price only if the state is inside the noise layer today); closed form, self-consistent: ((E/(M* theta)) sqrt(eps_now))^(2/3), E=40 GtCO2/yr (E_0), M*=7998 GtCO2 (calibration_ci.py), theta=9.677/yr, the vegetation-index time scale of Boulton et al. 2022, eps_now=0.66; theta declines as sqrt(eps) along the descent
\newcommand{\ampAmazonSlowSSPhigh}{1.0}   % closed form: F = (eps_now/eps_dyn)^(1/2) if eps_dyn < eps_now, else 1; the factor by which the leading law lifts the price between today's proximity and eps_dyn, the width of the window (the lift of the full price only if the state is inside the noise layer today); closed form, self-consistent: ((E/(M* theta)) sqrt(eps_now))^(2/3), E=82.7 GtCO2/yr (SSP3-7.0, highest rate before the median date, fig_forest_tc.py), M*=7998 GtCO2 (calibration_ci.py), theta=1/100/yr, assumption [sensitivity assumption, not a measurement: a recovery time of one century (yr)], eps_now=0.66; theta declines as sqrt(eps) along the descent
\newcommand{\ampAmazonSlowSSPmid}{1.1}   % closed form: F = (eps_now/eps_dyn)^(1/2) if eps_dyn < eps_now, else 1; the factor by which the leading law lifts the price between today's proximity and eps_dyn, the width of the window (the lift of the full price only if the state is inside the noise layer today); closed form, self-consistent: ((E/(M* theta)) sqrt(eps_now))^(2/3), E=44.3 GtCO2/yr (SSP2-4.5, highest rate before the median date, fig_forest_tc.py), M*=7998 GtCO2 (calibration_ci.py), theta=1/100/yr, assumption [sensitivity assumption, not a measurement: a recovery time of one century (yr)], eps_now=0.66; theta declines as sqrt(eps) along the descent
\newcommand{\ampAmazonSlow}{1.1}   % closed form: F = (eps_now/eps_dyn)^(1/2) if eps_dyn < eps_now, else 1; the factor by which the leading law lifts the price between today's proximity and eps_dyn, the width of the window (the lift of the full price only if the state is inside the noise layer today); closed form, self-consistent: ((E/(M* theta)) sqrt(eps_now))^(2/3), E=40 GtCO2/yr (E_0), M*=7998 GtCO2 (calibration_ci.py), theta=1/100/yr, assumption [sensitivity assumption, not a measurement: a recovery time of one century (yr)], eps_now=0.66; theta declines as sqrt(eps) along the descent
\newcommand{\ampWAIScenturySSPhigh}{1.0}   % closed form: F = (eps_now/eps_dyn)^(1/2) if eps_dyn < eps_now, else 1; the factor by which the leading law lifts the price between today's proximity and eps_dyn, the width of the window (the lift of the full price only if the state is inside the noise layer today); closed form, self-consistent: ((E/(M* theta)) sqrt(eps_now))^(2/3), E=59.8 GtCO2/yr (SSP3-7.0, highest rate before the median date, fig_forest_tc.py), M*=3895 GtCO2 (calibration_ci.py), theta=1/100/yr, slower modes (Robel et al. 2018), assumption [sensitivity assumption, not a measurement: a recovery time of one century (yr)], eps_now=0.2; theta declines as sqrt(eps) along the descent
\newcommand{\ampWAIScenturySSPmid}{1.0}   % closed form: F = (eps_now/eps_dyn)^(1/2) if eps_dyn < eps_now, else 1; the factor by which the leading law lifts the price between today's proximity and eps_dyn, the width of the window (the lift of the full price only if the state is inside the noise layer today); closed form, self-consistent: ((E/(M* theta)) sqrt(eps_now))^(2/3), E=44.3 GtCO2/yr (SSP2-4.5, highest rate before the median date, fig_forest_tc.py), M*=3895 GtCO2 (calibration_ci.py), theta=1/100/yr, slower modes (Robel et al. 2018), assumption [sensitivity assumption, not a measurement: a recovery time of one century (yr)], eps_now=0.2; theta declines as sqrt(eps) along the descent
\newcommand{\ampWAIScentury}{1.0}   % closed form: F = (eps_now/eps_dyn)^(1/2) if eps_dyn < eps_now, else 1; the factor by which the leading law lifts the price between today's proximity and eps_dyn, the width of the window (the lift of the full price only if the state is inside the noise layer today); closed form, self-consistent: ((E/(M* theta)) sqrt(eps_now))^(2/3), E=40 GtCO2/yr (E_0), M*=3895 GtCO2 (calibration_ci.py), theta=1/100/yr, slower modes (Robel et al. 2018), assumption [sensitivity assumption, not a measurement: a recovery time of one century (yr)], eps_now=0.2; theta declines as sqrt(eps) along the descent
\newcommand{\ampWAISfastSSPhigh}{1.1}   % closed form: F = (eps_now/eps_dyn)^(1/2) if eps_dyn < eps_now, else 1; the factor by which the leading law lifts the price between today's proximity and eps_dyn, the width of the window (the lift of the full price only if the state is inside the noise layer today); closed form, self-consistent: ((E/(M* theta)) sqrt(eps_now))^(2/3), E=59.8 GtCO2/yr (SSP3-7.0, highest rate before the median date, fig_forest_tc.py), M*=3895 GtCO2 (calibration_ci.py), theta=1/10/yr, flux relaxation, Hill et al. 2023 (relaxWAISHillLo), eps_now=0.2; theta declines as sqrt(eps) along the descent
\newcommand{\ampWAISfastSSPmid}{1.2}   % closed form: F = (eps_now/eps_dyn)^(1/2) if eps_dyn < eps_now, else 1; the factor by which the leading law lifts the price between today's proximity and eps_dyn, the width of the window (the lift of the full price only if the state is inside the noise layer today); closed form, self-consistent: ((E/(M* theta)) sqrt(eps_now))^(2/3), E=44.3 GtCO2/yr (SSP2-4.5, highest rate before the median date, fig_forest_tc.py), M*=3895 GtCO2 (calibration_ci.py), theta=1/10/yr, flux relaxation, Hill et al. 2023 (relaxWAISHillLo), eps_now=0.2; theta declines as sqrt(eps) along the descent
\newcommand{\ampWAISfast}{1.2}   % closed form: F = (eps_now/eps_dyn)^(1/2) if eps_dyn < eps_now, else 1; the factor by which the leading law lifts the price between today's proximity and eps_dyn, the width of the window (the lift of the full price only if the state is inside the noise layer today); closed form, self-consistent: ((E/(M* theta)) sqrt(eps_now))^(2/3), E=40 GtCO2/yr (E_0), M*=3895 GtCO2 (calibration_ci.py), theta=1/10/yr, flux relaxation, Hill et al. 2023 (relaxWAISHillLo), eps_now=0.2; theta declines as sqrt(eps) along the descent
\newcommand{\ampWAISslowSSPhigh}{1.0}   % closed form: F = (eps_now/eps_dyn)^(1/2) if eps_dyn < eps_now, else 1; the factor by which the leading law lifts the price between today's proximity and eps_dyn, the width of the window (the lift of the full price only if the state is inside the noise layer today); closed form, self-consistent: ((E/(M* theta)) sqrt(eps_now))^(2/3), E=59.8 GtCO2/yr (SSP3-7.0, highest rate before the median date, fig_forest_tc.py), M*=3895 GtCO2 (calibration_ci.py), theta=1/20/yr, flux relaxation, Hill et al. 2023 (relaxWAISHillHi), eps_now=0.2; theta declines as sqrt(eps) along the descent
\newcommand{\ampWAISslowSSPmid}{1.0}   % closed form: F = (eps_now/eps_dyn)^(1/2) if eps_dyn < eps_now, else 1; the factor by which the leading law lifts the price between today's proximity and eps_dyn, the width of the window (the lift of the full price only if the state is inside the noise layer today); closed form, self-consistent: ((E/(M* theta)) sqrt(eps_now))^(2/3), E=44.3 GtCO2/yr (SSP2-4.5, highest rate before the median date, fig_forest_tc.py), M*=3895 GtCO2 (calibration_ci.py), theta=1/20/yr, flux relaxation, Hill et al. 2023 (relaxWAISHillHi), eps_now=0.2; theta declines as sqrt(eps) along the descent
\newcommand{\ampWAISslow}{1.0}   % closed form: F = (eps_now/eps_dyn)^(1/2) if eps_dyn < eps_now, else 1; the factor by which the leading law lifts the price between today's proximity and eps_dyn, the width of the window (the lift of the full price only if the state is inside the noise layer today); closed form, self-consistent: ((E/(M* theta)) sqrt(eps_now))^(2/3), E=40 GtCO2/yr (E_0), M*=3895 GtCO2 (calibration_ci.py), theta=1/20/yr, flux relaxation, Hill et al. 2023 (relaxWAISHillHi), eps_now=0.2; theta declines as sqrt(eps) along the descent
\newcommand{\bhpCascadeFactor}{0.6116}   % bhp_bound.py — closed form of eq:ceiling_body under the floored hazard lambda_1 = h_1T max(0, T - T*), inputs from van2026three's replication package (Zenodo 20075290)
\newcommand{\bhpHoneRatio}{174}   % bhp_bound.py — closed form of eq:ceiling_body under the floored hazard lambda_1 = h_1T max(0, T - T*), inputs from van2026three's replication package (Zenodo 20075290)
\newcommand{\bhpHoneStar}{60.96\%\,{}^{\circ}\mathrm{C}^{-1}\,\mathrm{yr}^{-1}}   % bhp_bound.py — closed form of eq:ceiling_body under the floored hazard lambda_1 = h_1T max(0, T - T*), inputs from van2026three's replication package (Zenodo 20075290)
\newcommand{\bhpLamTwo}{0.08}   % bhp_bound.py — closed form of eq:ceiling_body under the floored hazard lambda_1 = h_1T max(0, T - T*), inputs from van2026three's replication package (Zenodo 20075290)
\newcommand{\bhpLambdaStar}{1.20}   % bhp_bound.py — closed form of eq:ceiling_body under the floored hazard lambda_1 = h_1T max(0, T - T*), inputs from van2026three's replication package (Zenodo 20075290)
\newcommand{\bhpLambda}{0.00689}   % bhp_bound.py — closed form of eq:ceiling_body under the floored hazard lambda_1 = h_1T max(0, T - T*), inputs from van2026three's replication package (Zenodo 20075290)
\newcommand{\bhpPriceRatio}{40.1}   % bhp_bound.py — closed form of eq:ceiling_body under the floored hazard lambda_1 = h_1T max(0, T - T*), inputs from van2026three's replication package (Zenodo 20075290)
\newcommand{\bhpRho}{0.0508}   % bhp_bound.py — closed form of eq:ceiling_body under the floored hazard lambda_1 = h_1T max(0, T - T*), inputs from van2026three's replication package (Zenodo 20075290)
\newcommand{\bhpShare}{2.49}   % bhp_bound.py — closed form of eq:ceiling_body under the floored hazard lambda_1 = h_1T max(0, T - T*), inputs from van2026three's replication package (Zenodo 20075290)
\newcommand{\bhpS}{0.0909}   % bhp_bound.py — closed form of eq:ceiling_body under the floored hazard lambda_1 = h_1T max(0, T - T*), inputs from van2026three's replication package (Zenodo 20075290)
\newcommand{\bhpWait}{2857}   % bhp_bound.py — closed form of eq:ceiling_body under the floored hazard lambda_1 = h_1T max(0, T - T*), inputs from van2026three's replication package (Zenodo 20075290)
\newcommand{\classFactorAtEpsDynAMOC}{6.1}   % closed form: eps_dyn^(-1/2) for the AMOC at E_0, eps_dyn=0.0272; the class factor tied at eps = 1, as \xoverClassFactorAtCentral  [requires network]
\newcommand{\emissionPathSSPmidZeroYear}{2250}   % fig_forest_tc.py via emission_paths.py — first year from which SSP2-4.5 CO2 emissions stay at zero in its published extension (Meinshausen et al. 2020); RCMIP v5.1.0 (Zenodo 4589756), Emissions|CO2, World, with the extensions of Meinshausen et al. 2020; extract in 00_data/emissions/
\newcommand{\emissionPathStartYear}{2024}   % fig_forest_tc.py via emission_paths.py — T0, the year the emission paths start from; 2024, the start year the paper adopts (E_0 and DT_0 carry no year of their own)
\newcommand{\epsAMOCnat}{0.70}   % closed form: (DT*-1.2)/DT*, DT*=4 (AM2022), DT_cur=1.2 (AR6)
\newcommand{\epsAmazonnat}{0.66}   % closed form: (DT*-1.2)/DT*, DT*=3.5 (AM2022), DT_cur=1.2 (AR6)
\newcommand{\epsDynAMOCSSPhigh}{0.044}   % closed form, self-consistent: ((E/(M* theta)) sqrt(eps_now))^(2/3), E=82.7 GtCO2/yr (SSP3-7.0, highest rate before the median date, fig_forest_tc.py), M*=9433 GtCO2 (calibration_ci.py), theta=0.7928/yr measured on the record (ews_calibration.py), eps_now=0.7; theta declines as sqrt(eps) along the descent  [requires network]
\newcommand{\epsDynAMOCSSPmid}{0.029}   % closed form, self-consistent: ((E/(M* theta)) sqrt(eps_now))^(2/3), E=44.3 GtCO2/yr (SSP2-4.5, highest rate before the median date, fig_forest_tc.py), M*=9433 GtCO2 (calibration_ci.py), theta=0.7928/yr measured on the record (ews_calibration.py), eps_now=0.7; theta declines as sqrt(eps) along the descent  [requires network]
\newcommand{\epsDynAMOC}{0.027}   % closed form, self-consistent: ((E/(M* theta)) sqrt(eps_now))^(2/3), E=40 GtCO2/yr (E_0), M*=9433 GtCO2 (calibration_ci.py), theta=0.7928/yr measured on the record (ews_calibration.py), eps_now=0.7; theta declines as sqrt(eps) along the descent  [requires network]
\newcommand{\epsDynAmazonBoultonSSPhigh}{0.0091}   % closed form, self-consistent: ((E/(M* theta)) sqrt(eps_now))^(2/3), E=82.7 GtCO2/yr (SSP3-7.0, highest rate before the median date, fig_forest_tc.py), M*=7998 GtCO2 (calibration_ci.py), theta=9.677/yr, the vegetation-index time scale of Boulton et al. 2022, eps_now=0.66; theta declines as sqrt(eps) along the descent
\newcommand{\epsDynAmazonBoultonSSPmid}{0.0060}   % closed form, self-consistent: ((E/(M* theta)) sqrt(eps_now))^(2/3), E=44.3 GtCO2/yr (SSP2-4.5, highest rate before the median date, fig_forest_tc.py), M*=7998 GtCO2 (calibration_ci.py), theta=9.677/yr, the vegetation-index time scale of Boulton et al. 2022, eps_now=0.66; theta declines as sqrt(eps) along the descent
\newcommand{\epsDynAmazonBoulton}{0.0056}   % closed form, self-consistent: ((E/(M* theta)) sqrt(eps_now))^(2/3), E=40 GtCO2/yr (E_0), M*=7998 GtCO2 (calibration_ci.py), theta=9.677/yr, the vegetation-index time scale of Boulton et al. 2022, eps_now=0.66; theta declines as sqrt(eps) along the descent
\newcommand{\epsDynAmazonSlowSSPhigh}{0.89}   % closed form, self-consistent: ((E/(M* theta)) sqrt(eps_now))^(2/3), E=82.7 GtCO2/yr (SSP3-7.0, highest rate before the median date, fig_forest_tc.py), M*=7998 GtCO2 (calibration_ci.py), theta=1/100/yr, assumption [sensitivity assumption, not a measurement: a recovery time of one century (yr)], eps_now=0.66; theta declines as sqrt(eps) along the descent
\newcommand{\epsDynAmazonSlowSSPmid}{0.59}   % closed form, self-consistent: ((E/(M* theta)) sqrt(eps_now))^(2/3), E=44.3 GtCO2/yr (SSP2-4.5, highest rate before the median date, fig_forest_tc.py), M*=7998 GtCO2 (calibration_ci.py), theta=1/100/yr, assumption [sensitivity assumption, not a measurement: a recovery time of one century (yr)], eps_now=0.66; theta declines as sqrt(eps) along the descent
\newcommand{\epsDynAmazonSlow}{0.55}   % closed form, self-consistent: ((E/(M* theta)) sqrt(eps_now))^(2/3), E=40 GtCO2/yr (E_0), M*=7998 GtCO2 (calibration_ci.py), theta=1/100/yr, assumption [sensitivity assumption, not a measurement: a recovery time of one century (yr)], eps_now=0.66; theta declines as sqrt(eps) along the descent
\newcommand{\epsDynWAIScenturySSPhigh}{0.78}   % closed form, self-consistent: ((E/(M* theta)) sqrt(eps_now))^(2/3), E=59.8 GtCO2/yr (SSP3-7.0, highest rate before the median date, fig_forest_tc.py), M*=3895 GtCO2 (calibration_ci.py), theta=1/100/yr, slower modes (Robel et al. 2018), assumption [sensitivity assumption, not a measurement: a recovery time of one century (yr)], eps_now=0.2; theta declines as sqrt(eps) along the descent
\newcommand{\epsDynWAIScenturySSPmid}{0.64}   % closed form, self-consistent: ((E/(M* theta)) sqrt(eps_now))^(2/3), E=44.3 GtCO2/yr (SSP2-4.5, highest rate before the median date, fig_forest_tc.py), M*=3895 GtCO2 (calibration_ci.py), theta=1/100/yr, slower modes (Robel et al. 2018), assumption [sensitivity assumption, not a measurement: a recovery time of one century (yr)], eps_now=0.2; theta declines as sqrt(eps) along the descent
\newcommand{\epsDynWAIScentury}{0.60}   % closed form, self-consistent: ((E/(M* theta)) sqrt(eps_now))^(2/3), E=40 GtCO2/yr (E_0), M*=3895 GtCO2 (calibration_ci.py), theta=1/100/yr, slower modes (Robel et al. 2018), assumption [sensitivity assumption, not a measurement: a recovery time of one century (yr)], eps_now=0.2; theta declines as sqrt(eps) along the descent
\newcommand{\epsDynWAISfastSSPhigh}{0.17}   % closed form, self-consistent: ((E/(M* theta)) sqrt(eps_now))^(2/3), E=59.8 GtCO2/yr (SSP3-7.0, highest rate before the median date, fig_forest_tc.py), M*=3895 GtCO2 (calibration_ci.py), theta=1/10/yr, flux relaxation, Hill et al. 2023 (relaxWAISHillLo), eps_now=0.2; theta declines as sqrt(eps) along the descent
\newcommand{\epsDynWAISfastSSPmid}{0.14}   % closed form, self-consistent: ((E/(M* theta)) sqrt(eps_now))^(2/3), E=44.3 GtCO2/yr (SSP2-4.5, highest rate before the median date, fig_forest_tc.py), M*=3895 GtCO2 (calibration_ci.py), theta=1/10/yr, flux relaxation, Hill et al. 2023 (relaxWAISHillLo), eps_now=0.2; theta declines as sqrt(eps) along the descent
\newcommand{\epsDynWAISfast}{0.13}   % closed form, self-consistent: ((E/(M* theta)) sqrt(eps_now))^(2/3), E=40 GtCO2/yr (E_0), M*=3895 GtCO2 (calibration_ci.py), theta=1/10/yr, flux relaxation, Hill et al. 2023 (relaxWAISHillLo), eps_now=0.2; theta declines as sqrt(eps) along the descent
\newcommand{\epsDynWAISslowSSPhigh}{0.27}   % closed form, self-consistent: ((E/(M* theta)) sqrt(eps_now))^(2/3), E=59.8 GtCO2/yr (SSP3-7.0, highest rate before the median date, fig_forest_tc.py), M*=3895 GtCO2 (calibration_ci.py), theta=1/20/yr, flux relaxation, Hill et al. 2023 (relaxWAISHillHi), eps_now=0.2; theta declines as sqrt(eps) along the descent
\newcommand{\epsDynWAISslowSSPmid}{0.22}   % closed form, self-consistent: ((E/(M* theta)) sqrt(eps_now))^(2/3), E=44.3 GtCO2/yr (SSP2-4.5, highest rate before the median date, fig_forest_tc.py), M*=3895 GtCO2 (calibration_ci.py), theta=1/20/yr, flux relaxation, Hill et al. 2023 (relaxWAISHillHi), eps_now=0.2; theta declines as sqrt(eps) along the descent
\newcommand{\epsDynWAISslow}{0.20}   % closed form, self-consistent: ((E/(M* theta)) sqrt(eps_now))^(2/3), E=40 GtCO2/yr (E_0), M*=3895 GtCO2 (calibration_ci.py), theta=1/20/yr, flux relaxation, Hill et al. 2023 (relaxWAISHillHi), eps_now=0.2; theta declines as sqrt(eps) along the descent
\newcommand{\epsRhoAMOC}{0.14}   % moving_budget.py — the price at fixed state when the budget is spent at E_0, against the frozen price (AMOC, eta*_now in {0.1, 0.25, 0.5, 1}); see the script's tables  [requires network]
\newcommand{\epsWAISnat}{0.20}   % closed form: (DT*-1.2)/DT*, DT*=1.5 (AM2022), DT_cur=1.2 (AR6)
\newcommand{\etaStarReversal}{1.00}   % noise_reversal_frontier.py — eta* at which dPi/dsigma changes sign for the fold, AT THE AMOC CALIBRATION: it is not a constant, moving with kappa = 2 rho/lambda, the kappa of the AMOC calibration printed in section 6 of the capture  [requires network]
\newcommand{\fracWAISalreadyPct}{25.2}   % fig_forest_tc.py — % of WAIS draws with DT* below present warming, i.e. mu_now = 0 and the deadline already passed (seed=42); these are excluded from the ordering sample
\newcommand{\fracWAISbeforeAmazonIncl}{96.4}   % fig_forest_tc.py — CONTROL on the exclusion: the same share with the already-crossed draws entered at t_c = 0 (seed=42)
\newcommand{\fracWAISbeforeAmazon}{95}   % fig_forest_tc.py — share of the draws with both bounds finite (t_c > 0) in which the ice sheet's bound is no later than the Amazon's (seed=42)
\newcommand{\kappaAMOC}{0.076}   % closed form: 2 rho / lambda, rho=0.03 (rhoBase), lambda=0.7928 (ews_calibration.py)  [requires network]
\newcommand{\kappaAmazonBoulton}{9.68}   % closed form: 12 / (1/kappa in months), 1/kappa=1.24 months [Boulton, Lenton and Boers 2022, Nature Climate Change 12:271-278, Methods p. 9 — 1/kappa of the VOD vegetation index, regional mean AR(1), months]; per year
\newcommand{\kappaReversalCeiling}{0.075}   % noise_reversal_frontier.py — the kappa below which the frontier lies ABOVE the window ceiling eta* = 1, so that the ceiling alone certifies the sign of dPi/dsigma  [requires network]
\newcommand{\movingDepartHi}{1.7}   % moving_budget.py — the price at fixed state when the budget is spent at E_0, against the frozen price (AMOC, eta*_now in {0.1, 0.25, 0.5, 1}); see the script's tables  [requires network]
\newcommand{\movingDepartLo}{0.2}   % moving_budget.py — the price at fixed state when the budget is spent at E_0, against the frozen price (AMOC, eta*_now in {0.1, 0.25, 0.5, 1}); see the script's tables  [requires network]
\newcommand{\movingErrNowPct}{11}   % moving_budget.py — the price at fixed state when the budget is spent at E_0, against the frozen price (AMOC, eta*_now in {0.1, 0.25, 0.5, 1}); see the script's tables  [requires network]
\newcommand{\movingPeakRatio}{1.8}   % moving_budget.py — the price at fixed state when the budget is spent at E_0, against the frozen price (AMOC, eta*_now in {0.1, 0.25, 0.5, 1}); see the script's tables  [requires network]
\newcommand{\movingRatioAtEpsRhoPctEtaOne}{8.5}   % moving_budget.py — the price at fixed state when the budget is spent at E_0, against the frozen price (AMOC, eta*_now in {0.1, 0.25, 0.5, 1}); see the script's tables  [requires network]
\newcommand{\movingRatioAtEpsRhoPct}{0.38}   % moving_budget.py — the price at fixed state when the budget is spent at E_0, against the frozen price (AMOC, eta*_now in {0.1, 0.25, 0.5, 1}); see the script's tables  [requires network]
\newcommand{\movingSlopeDeep}{-0.20}   % moving_budget.py — the price at fixed state when the budget is spent at E_0, against the frozen price (AMOC, eta*_now in {0.1, 0.25, 0.5, 1}); see the script's tables  [requires network]
\newcommand{\recBiasFailPct}{69}   % record_detectability.py — read from the record-detectability study outputs in 02_output/record_study (written by 01_code/05_record_study/, or copied from 00_data/record_study with RECORD_STUDY=cached)
\newcommand{\recBiasPassPct}{28}   % record_detectability.py — read from the record-detectability study outputs in 02_output/record_study (written by 01_code/05_record_study/, or copied from 00_data/record_study with RECORD_STUDY=cached)
\newcommand{\recClassPowerGap}{0.04}   % record_detectability.py — read from the record-detectability study outputs in 02_output/record_study (written by 01_code/05_record_study/, or copied from 00_data/record_study with RECORD_STUDY=cached)
\newcommand{\recDDConstSizeCentreA}{28}   % record_detectability.py — read from the record-detectability study outputs in 02_output/record_study (written by 01_code/05_record_study/, or copied from 00_data/record_study with RECORD_STUDY=cached)
\newcommand{\recDDConstSizeCentreB}{40}   % record_detectability.py — read from the record-detectability study outputs in 02_output/record_study (written by 01_code/05_record_study/, or copied from 00_data/record_study with RECORD_STUDY=cached)
\newcommand{\recDDConstSizeHi}{94}   % record_detectability.py — read from the record-detectability study outputs in 02_output/record_study (written by 01_code/05_record_study/, or copied from 00_data/record_study with RECORD_STUDY=cached)
\newcommand{\recDDConstSizeLo}{1.5}   % record_detectability.py — read from the record-detectability study outputs in 02_output/record_study (written by 01_code/05_record_study/, or copied from 00_data/record_study with RECORD_STUDY=cached)
\newcommand{\recDDConstSizeSigmaRising}{82}   % record_detectability.py — read from the record-detectability study outputs in 02_output/record_study (written by 01_code/05_record_study/, or copied from 00_data/record_study with RECORD_STUDY=cached)
\newcommand{\recDDFreeShareOfBoundHi}{81}   % record_detectability.py — read from the record-detectability study outputs in 02_output/record_study (written by 01_code/05_record_study/, or copied from 00_data/record_study with RECORD_STUDY=cached)
\newcommand{\recDDFreeShareOfBoundLo}{79}   % record_detectability.py — read from the record-detectability study outputs in 02_output/record_study (written by 01_code/05_record_study/, or copied from 00_data/record_study with RECORD_STUDY=cached)
\newcommand{\recDDFreeSizeHi}{10}   % record_detectability.py — read from the record-detectability study outputs in 02_output/record_study (written by 01_code/05_record_study/, or copied from 00_data/record_study with RECORD_STUDY=cached)
\newcommand{\recDDFreeSizeLo}{4}   % record_detectability.py — read from the record-detectability study outputs in 02_output/record_study (written by 01_code/05_record_study/, or copied from 00_data/record_study with RECORD_STUDY=cached)
\newcommand{\recHorizonFoldEightyFromNow}{128}   % record_detectability.py — read from the record-detectability study outputs in 02_output/record_study (written by 01_code/05_record_study/, or copied from 00_data/record_study with RECORD_STUDY=cached)
\newcommand{\recHorizonFoldEightyRecord}{274}   % record_detectability.py — read from the record-detectability study outputs in 02_output/record_study (written by 01_code/05_record_study/, or copied from 00_data/record_study with RECORD_STUDY=cached)
\newcommand{\recHorizonFoldEightySpentPct}{82}   % record_detectability.py — read from the record-detectability study outputs in 02_output/record_study (written by 01_code/05_record_study/, or copied from 00_data/record_study with RECORD_STUDY=cached)
\newcommand{\recHorizonFoldEightyToThreshold}{28.1}   % record_detectability.py — read from the record-detectability study outputs in 02_output/record_study (written by 01_code/05_record_study/, or copied from 00_data/record_study with RECORD_STUDY=cached)
\newcommand{\recHorizonFoldFiftyFixedRecord}{10.4}   % record_detectability.py — read from the record-detectability study outputs in 02_output/record_study (written by 01_code/05_record_study/, or copied from 00_data/record_study with RECORD_STUDY=cached)
\newcommand{\recHorizonFoldFiftyFromNow}{90}   % record_detectability.py — read from the record-detectability study outputs in 02_output/record_study (written by 01_code/05_record_study/, or copied from 00_data/record_study with RECORD_STUDY=cached)
\newcommand{\recHorizonFoldFiftyRecord}{236}   % record_detectability.py — read from the record-detectability study outputs in 02_output/record_study (written by 01_code/05_record_study/, or copied from 00_data/record_study with RECORD_STUDY=cached)
\newcommand{\recHorizonFoldFiftySpentPct}{58}   % record_detectability.py — read from the record-detectability study outputs in 02_output/record_study (written by 01_code/05_record_study/, or copied from 00_data/record_study with RECORD_STUDY=cached)
\newcommand{\recHorizonFoldFiftyToThreshold}{65.5}   % record_detectability.py — read from the record-detectability study outputs in 02_output/record_study (written by 01_code/05_record_study/, or copied from 00_data/record_study with RECORD_STUDY=cached)
\newcommand{\recHorizonTransEightyFromNow}{73}   % record_detectability.py — read from the record-detectability study outputs in 02_output/record_study (written by 01_code/05_record_study/, or copied from 00_data/record_study with RECORD_STUDY=cached)
\newcommand{\recHorizonTransEightyToThreshold}{83}   % record_detectability.py — read from the record-detectability study outputs in 02_output/record_study (written by 01_code/05_record_study/, or copied from 00_data/record_study with RECORD_STUDY=cached)
\newcommand{\recHorizonTransFiftyFromNow}{38}   % record_detectability.py — read from the record-detectability study outputs in 02_output/record_study (written by 01_code/05_record_study/, or copied from 00_data/record_study with RECORD_STUDY=cached)
\newcommand{\recHorizonTransFiftyToThreshold}{118}   % record_detectability.py — read from the record-detectability study outputs in 02_output/record_study (written by 01_code/05_record_study/, or copied from 00_data/record_study with RECORD_STUDY=cached)
\newcommand{\recJointOriginIn}{116}   % record_detectability.py — read from the record-detectability study outputs in 02_output/record_study (written by 01_code/05_record_study/, or copied from 00_data/record_study with RECORD_STUDY=cached)
\newcommand{\recJointRaysIn}{120}   % record_detectability.py — read from the record-detectability study outputs in 02_output/record_study (written by 01_code/05_record_study/, or copied from 00_data/record_study with RECORD_STUDY=cached)
\newcommand{\recJointSigChange}{0}   % record_detectability.py — read from the record-detectability study outputs in 02_output/record_study (written by 01_code/05_record_study/, or copied from 00_data/record_study with RECORD_STUDY=cached)
\newcommand{\recJointSpecs}{120}   % record_detectability.py — read from the record-detectability study outputs in 02_output/record_study (written by 01_code/05_record_study/, or copied from 00_data/record_study with RECORD_STUDY=cached)
\newcommand{\recJointThetaDecline}{22}   % record_detectability.py — read from the record-detectability study outputs in 02_output/record_study (written by 01_code/05_record_study/, or copied from 00_data/record_study with RECORD_STUDY=cached)
\newcommand{\recNPFoldNear}{0.68}   % record_detectability.py — read from the record-detectability study outputs in 02_output/record_study (written by 01_code/05_record_study/, or copied from 00_data/record_study with RECORD_STUDY=cached)
\newcommand{\recNPFoldNowGridHi}{0.13}   % record_detectability.py — read from the record-detectability study outputs in 02_output/record_study (written by 01_code/05_record_study/, or copied from 00_data/record_study with RECORD_STUDY=cached)
\newcommand{\recNPFoldNowGridLo}{0.11}   % record_detectability.py — read from the record-detectability study outputs in 02_output/record_study (written by 01_code/05_record_study/, or copied from 00_data/record_study with RECORD_STUDY=cached)
\newcommand{\recNPFoldNow}{0.12}   % record_detectability.py — read from the record-detectability study outputs in 02_output/record_study (written by 01_code/05_record_study/, or copied from 00_data/record_study with RECORD_STUDY=cached)
\newcommand{\recNPTransNow}{0.24}   % record_detectability.py — read from the record-detectability study outputs in 02_output/record_study (written by 01_code/05_record_study/, or copied from 00_data/record_study with RECORD_STUDY=cached)
\newcommand{\recNPtcNear}{0.5}   % record_detectability.py — read from the record-detectability study outputs in 02_output/record_study (written by 01_code/05_record_study/, or copied from 00_data/record_study with RECORD_STUDY=cached)
\newcommand{\recNPtcNow}{156}   % record_detectability.py — read from the record-detectability study outputs in 02_output/record_study (written by 01_code/05_record_study/, or copied from 00_data/record_study with RECORD_STUDY=cached)
\newcommand{\recRecordYears}{146}   % record_detectability.py — read from the record-detectability study outputs in 02_output/record_study (written by 01_code/05_record_study/, or copied from 00_data/record_study with RECORD_STUDY=cached)
\newcommand{\recSEExponent}{3/2}   % record_detectability.py — read from the record-detectability study outputs in 02_output/record_study (written by 01_code/05_record_study/, or copied from 00_data/record_study with RECORD_STUDY=cached)
\newcommand{\recSigmaTrendCentre}{0.46}   % record_detectability.py — read from the record-detectability study outputs in 02_output/record_study (written by 01_code/05_record_study/, or copied from 00_data/record_study with RECORD_STUDY=cached)
\newcommand{\recSigmaTrendHi}{1.06}   % record_detectability.py — read from the record-detectability study outputs in 02_output/record_study (written by 01_code/05_record_study/, or copied from 00_data/record_study with RECORD_STUDY=cached)
\newcommand{\recSigmaTrendLo}{-0.15}   % record_detectability.py — read from the record-detectability study outputs in 02_output/record_study (written by 01_code/05_record_study/, or copied from 00_data/record_study with RECORD_STUDY=cached)
\newcommand{\recSlopeLimitCoef}{3}   % record_detectability.py — read from the record-detectability study outputs in 02_output/record_study (written by 01_code/05_record_study/, or copied from 00_data/record_study with RECORD_STUDY=cached)
\newcommand{\recThetaWFail}{4.5}   % record_detectability.py — read from the record-detectability study outputs in 02_output/record_study (written by 01_code/05_record_study/, or copied from 00_data/record_study with RECORD_STUDY=cached)
\newcommand{\recThetaWFloor}{10}   % record_detectability.py — read from the record-detectability study outputs in 02_output/record_study (written by 01_code/05_record_study/, or copied from 00_data/record_study with RECORD_STUDY=cached)
\newcommand{\recThetaWMinValid}{8.9}   % record_detectability.py — read from the record-detectability study outputs in 02_output/record_study (written by 01_code/05_record_study/, or copied from 00_data/record_study with RECORD_STUDY=cached)
\newcommand{\recThetaWPass}{11.4}   % record_detectability.py — read from the record-detectability study outputs in 02_output/record_study (written by 01_code/05_record_study/, or copied from 00_data/record_study with RECORD_STUDY=cached)
\newcommand{\recTrendLengthFifty}{349}   % record_detectability.py — read from the record-detectability study outputs in 02_output/record_study (written by 01_code/05_record_study/, or copied from 00_data/record_study with RECORD_STUDY=cached)
\newcommand{\recWindowCells}{16}   % record_detectability.py — read from the record-detectability study outputs in 02_output/record_study (written by 01_code/05_record_study/, or copied from 00_data/record_study with RECORD_STUDY=cached)
\newcommand{\recWindowOptCells}{15}   % record_detectability.py — read from the record-detectability study outputs in 02_output/record_study (written by 01_code/05_record_study/, or copied from 00_data/record_study with RECORD_STUDY=cached)
\newcommand{\recWindowOptSharePct}{41}   % record_detectability.py — read from the record-detectability study outputs in 02_output/record_study (written by 01_code/05_record_study/, or copied from 00_data/record_study with RECORD_STUDY=cached)
\newcommand{\recWindowOptYears}{60}   % record_detectability.py — read from the record-detectability study outputs in 02_output/record_study (written by 01_code/05_record_study/, or copied from 00_data/record_study with RECORD_STUDY=cached)
\newcommand{\recWindowYears}{40}   % record_detectability.py — read from the record-detectability study outputs in 02_output/record_study (written by 01_code/05_record_study/, or copied from 00_data/record_study with RECORD_STUDY=cached)
\newcommand{\sccReducedCoefAMOC}{0.033}   % reduced_form_scc.py — the reduced-form price at the median damage prior (lambda, Var from ews_calibration.py)  [requires network]
\newcommand{\sdAMOCrecord}{0.30}   % reduced_form_scc.py — the reduced-form price at the median damage prior (lambda, Var from ews_calibration.py)  [requires network]
\newcommand{\sigmaAMOCPninety}{0.54}   % closed form: lognormal(median=0.3725, cv=0.3) percentiles [prior CV on AMOC sigma (= sigmaPriorPct), for P10/P90]  [requires network]
\newcommand{\sigmaAMOCPten}{0.26}   % closed form: lognormal(median=0.3725, cv=0.3) percentiles [prior CV on AMOC sigma (= sigmaPriorPct), for P10/P90]  [requires network]
\newcommand{\sigmaAMOChat}{0.372}   % ews_calibration.py — sigma_clim detrended  [requires network]
\newcommand{\sigmaAggregationFactorTheta}{1.27}   % closed form: the same factor through the recovery rate alone, sqrt(lambda/-ln AC1), lambda=0.7928, -ln AC1=0.4933 (ews_calibration.py)  [requires network]
\newcommand{\sigmaAggregationFactor}{1.43}   % closed form: sigma-hat (annual-mean inversion) / sigma-hat (point-sample reading) = sqrt((lambda/-ln AC1) x Var(X_t)/Var(annual means)), lambda=0.7928, -ln AC1=0.4933, variance ratio=1.2808 (ews_calibration.py)  [requires network]
\newcommand{\tcAMOCPten}{47}   % fig_forest_tc.py — P10 of the deadline bound mu_now/E0 over the draws with t_c > 0, the left end of the fig:forest_tc bar (seed=42); same sample as the median, no floor
\newcommand{\tcAMOCSSPhigh}{95}   % fig_forest_tc.py — median of the deadline bound on SSP3-7.0: years from T0 until emissions summed from T0 reach mu_now, same draws and sample as \tcAMOC (seed=42); unreached draws count as +inf, a centile among them gets no macro
\newcommand{\tcAMOCSSPmidNotReachedPct}{75.0}   % fig_forest_tc.py — % of the draws with t_c > 0 that SSP2-4.5 never brings to mu_now before \emissionPathLastYear (seed=42)
\newcommand{\tcAMOC}{164}   % eps_inf_ci.py — deadline bound t_c <= (M*/E0) eps_0 = mu_now/E0 (median calibration, seed=42); free of the backstop cost and of the margin
\newcommand{\tcAmazonPten}{54}   % fig_forest_tc.py — P10 of the deadline bound mu_now/E0 over the draws with t_c > 0, the left end of the fig:forest_tc bar (seed=42); same sample as the median, no floor
\newcommand{\tcAmazonSSPhighNotReachedPct}{14.1}   % fig_forest_tc.py — % of the draws with t_c > 0 that SSP3-7.0 never brings to mu_now before \emissionPathLastYear (seed=42)
\newcommand{\tcAmazonSSPhighPten}{37}   % fig_forest_tc.py — P10 of the deadline bound on SSP3-7.0: years from T0 until emissions summed from T0 reach mu_now, same draws and sample as \tcAmazon (seed=42); unreached draws count as +inf, a centile among them gets no macro
\newcommand{\tcAmazonSSPhigh}{78}   % fig_forest_tc.py — median of the deadline bound on SSP3-7.0: years from T0 until emissions summed from T0 reach mu_now, same draws and sample as \tcAmazon (seed=42); unreached draws count as +inf, a centile among them gets no macro
\newcommand{\tcAmazonSSPmidNotReachedPct}{71.7}   % fig_forest_tc.py — % of the draws with t_c > 0 that SSP2-4.5 never brings to mu_now before \emissionPathLastYear (seed=42)
\newcommand{\tcAmazonSSPmidPten}{53}   % fig_forest_tc.py — P10 of the deadline bound on SSP2-4.5: years from T0 until emissions summed from T0 reach mu_now, same draws and sample as \tcAmazon (seed=42); unreached draws count as +inf, a centile among them gets no macro
\newcommand{\tcAmazon}{131}   % eps_inf_ci.py — deadline bound t_c <= (M*/E0) eps_0 = mu_now/E0 (median calibration, seed=42); free of the backstop cost and of the margin
\newcommand{\tcWAISPten}{5}   % fig_forest_tc.py — P10 of the deadline bound mu_now/E0 over the draws with t_c > 0, the left end of the fig:forest_tc bar (seed=42); same sample as the median, no floor
\newcommand{\tcWAISSSPhighNotReachedPct}{1.2}   % fig_forest_tc.py — % of the draws with t_c > 0 that SSP3-7.0 never brings to mu_now before \emissionPathLastYear (seed=42)
\newcommand{\tcWAISSSPhighPninety}{57}   % fig_forest_tc.py — P90 of the deadline bound on SSP3-7.0: years from T0 until emissions summed from T0 reach mu_now, same draws and sample as \tcWAIS (seed=42); unreached draws count as +inf, a centile among them gets no macro
\newcommand{\tcWAISSSPhighPten}{4}   % fig_forest_tc.py — P10 of the deadline bound on SSP3-7.0: years from T0 until emissions summed from T0 reach mu_now, same draws and sample as \tcWAIS (seed=42); unreached draws count as +inf, a centile among them gets no macro
\newcommand{\tcWAISSSPhigh}{20}   % fig_forest_tc.py — median of the deadline bound on SSP3-7.0: years from T0 until emissions summed from T0 reach mu_now, same draws and sample as \tcWAIS (seed=42); unreached draws count as +inf, a centile among them gets no macro
\newcommand{\tcWAISSSPmidNotReachedPct}{10.0}   % fig_forest_tc.py — % of the draws with t_c > 0 that SSP2-4.5 never brings to mu_now before \emissionPathLastYear (seed=42)
\newcommand{\tcWAISSSPmidPninety}{215}   % fig_forest_tc.py — P90 of the deadline bound on SSP2-4.5: years from T0 until emissions summed from T0 reach mu_now, same draws and sample as \tcWAIS (seed=42); unreached draws count as +inf, a centile among them gets no macro
\newcommand{\tcWAISSSPmidPten}{5}   % fig_forest_tc.py — P10 of the deadline bound on SSP2-4.5: years from T0 until emissions summed from T0 reach mu_now, same draws and sample as \tcWAIS (seed=42); unreached draws count as +inf, a centile among them gets no macro
\newcommand{\tcWAISSSPmid}{25}   % fig_forest_tc.py — median of the deadline bound on SSP2-4.5: years from T0 until emissions summed from T0 reach mu_now, same draws and sample as \tcWAIS (seed=42); unreached draws count as +inf, a centile among them gets no macro
\newcommand{\tcWAIS}{27}   % eps_inf_ci.py — deadline bound t_c <= (M*/E0) eps_0 = mu_now/E0 (median calibration, seed=42); free of the backstop cost and of the margin
\newcommand{\thetaAMOChat}{0.79}   % ews_calibration.py — AR(1) restoring rate  [requires network]
\newcommand{\valAmocRhoTilde}{0.0076}   % validation_bvp.py — app:solver, BVP validation (solve_bvp tol 1e-8, domain >= 8 noise scales); see the script's printed tables  [requires network]
\newcommand{\valBtermRatioAmoc}{2.6}   % validation_bvp.py — app:solver, BVP validation (solve_bvp tol 1e-8, domain >= 8 noise scales); see the script's printed tables  [requires network]
\newcommand{\valCollapseDevPctA}{0.0028}   % validation_bvp.py — app:solver, BVP validation (solve_bvp tol 1e-8, domain >= 8 noise scales); see the script's printed tables
\newcommand{\valCollapseDevPctB}{0.0025}   % validation_bvp.py — app:solver, BVP validation (solve_bvp tol 1e-8, domain >= 8 noise scales); see the script's printed tables
\newcommand{\valCollapseDevPctC}{0.0038}   % validation_bvp.py — app:solver, BVP validation (solve_bvp tol 1e-8, domain >= 8 noise scales); see the script's printed tables
\newcommand{\valCollapseGapPct}{0.0064}   % validation_bvp.py — app:solver, BVP validation (solve_bvp tol 1e-8, domain >= 8 noise scales); see the script's printed tables
\newcommand{\valCollapseRhoTildeA}{0.05}   % validation_bvp.py — app:solver, BVP validation (solve_bvp tol 1e-8, domain >= 8 noise scales); see the script's printed tables
\newcommand{\valCollapseRhoTildeB}{0.2}   % validation_bvp.py — app:solver, BVP validation (solve_bvp tol 1e-8, domain >= 8 noise scales); see the script's printed tables
\newcommand{\valCollapseRhoTildeC}{1}   % validation_bvp.py — app:solver, BVP validation (solve_bvp tol 1e-8, domain >= 8 noise scales); see the script's printed tables
\newcommand{\valFoldSlopeFar}{-0.492}   % validation_bvp.py — app:solver, BVP validation (solve_bvp tol 1e-8, domain >= 8 noise scales); see the script's printed tables
\newcommand{\valLawRatioAmoc}{1.2}   % validation_bvp.py — app:solver, BVP validation (solve_bvp tol 1e-8, domain >= 8 noise scales); see the script's printed tables  [requires network]
\newcommand{\valLawRatioSmallRhoA}{1.1}   % validation_bvp.py — app:solver, BVP validation (solve_bvp tol 1e-8, domain >= 8 noise scales); see the script's printed tables
\newcommand{\valLawRatioSmallRhoB}{1.9}   % validation_bvp.py — app:solver, BVP validation (solve_bvp tol 1e-8, domain >= 8 noise scales); see the script's printed tables
\newcommand{\valLeadEpsNarrow}{0.01}   % validation_bvp.py — app:solver, BVP validation (solve_bvp tol 1e-8, domain >= 8 noise scales); see the script's printed tables
\newcommand{\valLeadEpsWide}{0.05}   % validation_bvp.py — app:solver, BVP validation (solve_bvp tol 1e-8, domain >= 8 noise scales); see the script's printed tables
\newcommand{\valLeadErrMaxPctNarrow}{26}   % validation_bvp.py — app:solver, BVP validation (solve_bvp tol 1e-8, domain >= 8 noise scales); see the script's printed tables
\newcommand{\valLeadErrMaxPctWide}{94}   % validation_bvp.py — app:solver, BVP validation (solve_bvp tol 1e-8, domain >= 8 noise scales); see the script's printed tables
\newcommand{\valLeadErrSlopeHi}{0.63}   % validation_bvp.py — app:solver, BVP validation (solve_bvp tol 1e-8, domain >= 8 noise scales); see the script's printed tables
\newcommand{\valLeadErrSlopeLo}{0.61}   % validation_bvp.py — app:solver, BVP validation (solve_bvp tol 1e-8, domain >= 8 noise scales); see the script's printed tables
\newcommand{\valLeadErrSlope}{0.62}   % validation_bvp.py — app:solver, BVP validation (solve_bvp tol 1e-8, domain >= 8 noise scales); see the script's printed tables
\newcommand{\valLeadEtaMax}{0.32}   % validation_bvp.py — app:solver, BVP validation (solve_bvp tol 1e-8, domain >= 8 noise scales); see the script's printed tables
\newcommand{\valPitchSlopeFix}{0.50}   % validation_bvp.py — app:solver, BVP validation (solve_bvp tol 1e-8, domain >= 8 noise scales); see the script's printed tables
\newcommand{\valSensFEtaHi}{1}   % validation_bvp.py — app:solver, BVP validation (solve_bvp tol 1e-8, domain >= 8 noise scales); see the script's printed tables  [requires network]
\newcommand{\valSensFEtaLo}{0.1}   % validation_bvp.py — app:solver, BVP validation (solve_bvp tol 1e-8, domain >= 8 noise scales); see the script's printed tables  [requires network]
\newcommand{\valSensFShortPctHi}{87}   % validation_bvp.py — app:solver, BVP validation (solve_bvp tol 1e-8, domain >= 8 noise scales); see the script's printed tables  [requires network]
\newcommand{\valSensFShortPctLo}{11}   % validation_bvp.py — app:solver, BVP validation (solve_bvp tol 1e-8, domain >= 8 noise scales); see the script's printed tables  [requires network]
\newcommand{\valSmallRhoEtaA}{0.1}   % validation_bvp.py — app:solver, BVP validation (solve_bvp tol 1e-8, domain >= 8 noise scales); see the script's printed tables
\newcommand{\valSmallRhoEtaB}{0.3}   % validation_bvp.py — app:solver, BVP validation (solve_bvp tol 1e-8, domain >= 8 noise scales); see the script's printed tables
\newcommand{\valSmallRho}{0.02}   % validation_bvp.py — app:solver, BVP validation (solve_bvp tol 1e-8, domain >= 8 noise scales); see the script's printed tables
\newcommand{\valTcSlopeFix}{1.00}   % validation_bvp.py — app:solver, BVP validation (solve_bvp tol 1e-8, domain >= 8 noise scales); see the script's printed tables
\newcommand{\xoverClassFactorAtLow}{2.65}   % crossover.py — the power law SCC ~ eps^(p-1) read at p=1/2 and p=1 on the AMOC threshold range (Armstrong McKay et al. 2022, Table 1) and the current warming (Forster et al. 2024); the class factor ties the two prefactors at eps = 1
\newcommand{\xoverEpsHigh}{0.850}   % crossover.py — the power law SCC ~ eps^(p-1) read at p=1/2 and p=1 on the AMOC threshold range (Armstrong McKay et al. 2022, Table 1) and the current warming (Forster et al. 2024); the class factor ties the two prefactors at eps = 1
\newcommand{\xoverEpsLow}{0.143}   % crossover.py — the power law SCC ~ eps^(p-1) read at p=1/2 and p=1 on the AMOC threshold range (Armstrong McKay et al. 2022, Table 1) and the current warming (Forster et al. 2024); the class factor ties the two prefactors at eps = 1
\newcommand{\xoverEpsPct}{2.9}   % crossover.py — the power law SCC ~ eps^(p-1) read at p=1/2 and p=1 on the AMOC threshold range (Armstrong McKay et al. 2022, Table 1) and the current warming (Forster et al. 2024); the class factor ties the two prefactors at eps = 1
\newcommand{\xoverEps}{0.029}   % crossover.py — the power law SCC ~ eps^(p-1) read at p=1/2 and p=1 on the AMOC threshold range (Armstrong McKay et al. 2022, Table 1) and the current warming (Forster et al. 2024); the class factor ties the two prefactors at eps = 1
\newcommand{\xoverRangeFactor}{5.83}   % crossover.py — the power law SCC ~ eps^(p-1) read at p=1/2 and p=1 on the AMOC threshold range (Armstrong McKay et al. 2022, Table 1) and the current warming (Forster et al. 2024); the class factor ties the two prefactors at eps = 1
